% Springer Lecture Notes in Computer Science, on the llncs class that comes with TeX Live.
% Guidelines: https://www.springer.com/gp/computer-science/lncs/conference-proceedings-guidelines
\documentclass[runningheads]{llncs}
\usepackage[T1]{fontenc}
\usepackage{graphicx}
\usepackage{amsmath,amssymb}

\begin{document}

\title{Paper Title}
\author{First Author\inst{1}\orcidID{0000-0000-0000-0000} \and
Second Author\inst{2}}
\authorrunning{F. Author and S. Author}
\institute{Organization, City, Country\\
\email{email@example.com} \and
Organization, City, Country\\
\email{email@example.com}}

\maketitle

\begin{abstract}
A short summary of the problem, the approach and the results, in 150--250 words.
\keywords{First keyword \and Second keyword \and Third keyword.}
\end{abstract}

\section{Introduction}
Say what the problem is and why it matters, then what this paper does~\cite{lamport1994}.

\section{Method}
\begin{equation}
  y = f(x) + \epsilon
  \label{eq:model}
\end{equation}

\begin{figure}
  \centering
  \includegraphics[width=0.8\textwidth]{example-image}
  \caption{A figure.}
  \label{fig:example}
\end{figure}

\section{Experiments}
Results go here; see Fig.~\ref{fig:example}.

\section{Conclusion}
What was shown, and what comes next.

\begin{credits}
\subsubsection{\ackname} Funding and people to thank.

\subsubsection{\discintname}
The authors have no competing interests to declare that are relevant to the content of this article.
\end{credits}

\bibliographystyle{splncs04}
\bibliography{references}

\end{document}
